% Generated from validated Article Draft Result. Do not hand-edit; revise the structured draft instead.
\section{Multimodal — 「生成できる」から、リアルタイムに行き来できるへ}
\noindent\textbf{7月のmultimodal更新は、画像・音声・動画を個別に生成する競争だけではなかった。MiniMax H3は複数modalitiesをまたぐ生成とreferenceを掲げ、GPT-Liveは同時に聞き話すinteractionを前面に出し、Alibaba側ではQwen-Audio 3.0とWan 2.7のavailabilityが更新された。}\autocite{src-3fb14968493d,src-d8e13386974c,src-eac8dfbeb78c}

\subsection{MiniMax H3: multimodal generationを一つのmodel storyにまとめる}

MiniMaxは7月末にH3を発表し、image、video、audio generationとcross-modal context understandingを一つのmultimodal generative modelとして扱う方向を示した。発表では最大2K video output、native stereo sound、image・audio・videoをまたぐmultimodal referenceも説明されている。ここで重要なのは個別modalの最高点ではなく、異なる入力と出力を一つのinteractionへ接続しようとしている点である。\autocite{src-d8e13386974c}


\begin{claimboundary}
H3発表時点でMiniMaxはweightsを『今後数日で公開する予定』と説明していた。したがって7月31日のlaunchを、すでに完了したopen-weight releaseとして記述してはいけない。generation qualityもvendor demonstrationの範囲に留まる。\autocite{src-d8e13386974c}
\end{claimboundary}

\subsection{GPT-Live: multimodalをリアルタイムinteractionへ寄せる}

OpenAIは7月8日にGPT-Liveを、simultaneous listening and speakingを想定したreal-time modelとして発表した。公式説明では、live interaction loopの背後でより深い処理をfrontier modelへ委譲し、launch時にはGPT-5.5を利用する構成が示されている。これは『一つのmodelがすべてを同期的に処理する』構図ではなく、低遅延interactionと深い推論を役割分担するarchitectureとして読める。\autocite{src-3fb14968493d}


\begin{claimboundary}
GPT-Liveのavailability、latency、qualityに関する記述はOpenAIの一次情報に基づく。今回のEvidenceは独立したreal-time quality評価を確立するものではないため、interaction設計の存在と品質優位性を分離して扱う。\autocite{src-3fb14968493d}
\end{claimboundary}

\subsection{Qwen-AudioとWan: availabilityの積み重ねがecosystemを作る}

Alibaba Cloud Model Studioのmodel lifecycleでは、7月1日にWan 2.7 reference-to-video snapshot、7月14日にQwen-Audio 3.0 TTS modelsのinternational availabilityが記録されている。Qwen-Audio 3.0についてAlibabaはFlash variantを低遅延real-time interaction向け、Plus variantを高品質なprofessional scenario向けと位置付ける。月次で見ると、multimodal競争は大規模launchだけでなく、region別availabilityと用途別variantの整備でも進行している。\autocite{src-eac8dfbeb78c}


\begin{claimboundary}
Alibaba Model Studioのavailabilityはregion-specificであり、Flash/Plusのqualityやlatencyの位置付けはvendor claimである。利用可能になった事実と、どちらがどの程度優れているかという評価を混同しない。\autocite{src-eac8dfbeb78c}
\end{claimboundary}

\subsection{7月に見えた共通方向}

H3はcross-modal referenceを含む統合生成、GPT-Liveは同時listening/speakingとfrontier-model delegation、Qwen-Audio/Wanは音声・動画modelのavailability拡大を示した。3者は同じproductではないが、multimodalを『画像も音声も扱える』というfeature listから、時間軸を持つinteraction、複数modal間のreference、実際に呼び出せるserviceへ移している。今後の比較ではdemo qualityだけでなく、latency、streaming、reference consistency、region availability、weights公開状態を分けて追う必要がある。\autocite{src-3fb14968493d,src-d8e13386974c,src-eac8dfbeb78c}
